\textbf{Phần 1. Trắc nghiệm lựa chọn}
\setcounterpageref{ex}{0}

\begin{ex} %[5D10V60]
Một ô tô dự định đi từ tỉnh A đến tỉnh B cách nhau $180\text{ km}$ với vận tốc $45\text{ km/h}$. Nhưng sau khi đi được $2$ giờ, xe tăng vận tốc thêm $5\text{ km/h}$ trên quãng đường còn lại. Hỏi ô tô đến B sớm hơn dự định bao nhiêu phút?
\choice
{$10\text{ phút}$}
{\True $12\text{ phút}$}
{$15\text{ phút}$}
{$20\text{ phút}$}
\loigiai{
Thời gian dự định đi hết quãng đường là:
$$180 : 45 = 4\text{ (giờ)}.$$
Trong $2$ giờ đầu, ô tô đi được quãng đường là:
$$45 \times 2 = 90\text{ (km)}.$$
Quãng đường còn lại dài là:
$$180 - 90 = 90\text{ (km)}.$$
Vận tốc mới trên quãng đường còn lại là:
$$45 + 5 = 50\text{ (km/h)}.$$
Thời gian đi quãng đường còn lại là:
$$90 : 50 = 1{,}8\text{ (giờ)} = 1\text{ giờ } 48\text{ phút}.$$
Tổng thời gian đi thực tế là:
$$2\text{ giờ} + 1\text{ giờ } 48\text{ phút} = 3\text{ giờ } 48\text{ phút}.$$
Ô tô đến B sớm hơn dự định là:
$$4\text{ giờ} - 3\text{ giờ } 48\text{ phút} = 12\text{ (phút)}.$$
\textbf{Đáp án của câu hỏi là B}
}
\end{ex}

\begin{ex} %[5D10V60]
Một xe khách đi từ bến xe lúc $6\text{ giờ } 30\text{ phút}$ với vận tốc $48\text{ km/h}$. Dọc đường xe nghỉ trả khách mất $20\text{ phút}$. Hỏi xe khách tới điểm đến cách bến xe $120\text{ km}$ vào thời điểm nào?
\choice
{$9\text{ giờ } 10\text{ phút}$}
{\True $9\text{ giờ } 20\text{ phút}$}
{$9\text{ giờ } 30\text{ phút}$}
{$9\text{ giờ } 40\text{ phút}$}
\loigiai{
Thời gian xe chạy trên đường là:
$$120 : 48 = 2{,}5\text{ (giờ)} = 2\text{ giờ } 30\text{ phút}.$$
Tổng thời gian từ lúc xuất phát đến lúc tới nơi là:
$$2\text{ giờ } 30\text{ phút} + 20\text{ phút} = 2\text{ giờ } 50\text{ phút}.$$
Xe khách tới nơi lúc:
$$6\text{ giờ } 30\text{ phút} + 2\text{ giờ } 50\text{ phút} = 8\text{ giờ } 80\text{ phút} = 9\text{ giờ } 20\text{ phút}.$$
\textbf{Đáp án của câu hỏi là B}
}
\end{ex}

\begin{ex} %[5D10V60]
Một đoàn tàu hỏa dài $160\text{ m}$ chạy qua một đường hầm dài $840\text{ m}$ với vận tốc $45\text{ km/h}$. Thời gian từ lúc đầu tàu bắt đầu vào hầm đến lúc toa cuối cùng ra khỏi hầm là
\choice
{$1\text{ phút } 15\text{ giây}$}
{\True $1\text{ phút } 20\text{ giây}$}
{$1\text{ phút } 30\text{ giây}$}
{$1\text{ phút } 10\text{ giây}$}
\loigiai{
Quãng đường đoàn tàu phải đi là:
$$840 + 160 = 1\,000\text{ (m)} = 1\text{ km}.$$
Đổi: $45\text{ km/h} = \dfrac{45}{3{,}6}\text{ m/s} = 12{,}5\text{ m/s}$.\\
Thời gian đoàn tàu chui qua hầm là:
$$1\,000 : 12{,}5 = 80\text{ (giây)} = 1\text{ phút } 20\text{ giây}.$$
\textbf{Đáp án của câu hỏi là B}
}
\end{ex}

\begin{ex} %[5D10V60]
Một người đi xe đạp với vận tốc $12\text{ km/h}$ trong $2$ giờ $15$ phút. Sau đó người đó đi tiếp bằng xe buýt trong $45$ phút với vận tốc $40\text{ km/h}$. Tổng quãng đường người đó đã đi là
\choice
{$54\text{ km}$}
{\True $57\text{ km}$}
{$60\text{ km}$}
{$62\text{ km}$}
\loigiai{
Đổi: $2\text{ giờ } 15\text{ phút} = 2{,}25\text{ giờ}$; $45\text{ phút} = 0{,}75\text{ giờ}$.\\
Quãng đường đi bằng xe đạp là:
$$12 \times 2{,}25 = 27\text{ (km)}.$$
Quãng đường đi bằng xe buýt là:
$$40 \times 0{,}75 = 30\text{ (km)}.$$
Tổng quãng đường đã đi là:
$$27 + 30 = 57\text{ (km)}.$$
\textbf{Đáp án của câu hỏi là B}
}
\end{ex}

\begin{ex} %[5D10V60]
Bác Nam đi xe máy từ quê lên thành phố với vận tốc $40\text{ km/h}$ thì hết $3$ giờ $30$ phút. Nếu bác Nam đi bằng xe ô tô với vận tốc $56\text{ km/h}$ thì mất bao nhiêu thời gian?
\choice
{$2\text{ giờ } 15\text{ phút}$}
{\True $2\text{ giờ } 30\text{ phút}$}
{$2\text{ giờ } 45\text{ phút}$}
{$2\text{ giờ } 20\text{ phút}$}
\loigiai{
Đổi: $3\text{ giờ } 30\text{ phút} = 3{,}5\text{ giờ}$.\\
Quãng đường từ quê lên thành phố là:
$$40 \times 3{,}5 = 140\text{ (km)}.$$
Thời gian đi bằng ô tô là:
$$140 : 56 = 2{,}5\text{ (giờ)} = 2\text{ giờ } 30\text{ phút}.$$
\textbf{Đáp án của câu hỏi là B}
}
\end{ex}

\begin{ex} %[5D10V60]
Một người chạy thể dục với vận tốc không đổi $250\text{ m/phút}$. Trong $1$ giờ $12$ phút, người đó chạy được quãng đường là
\choice
{$16\text{ km}$}
{\True $18\text{ km}$}
{$20\text{ km}$}
{$15\text{ km}$}
\loigiai{
Đổi: $1\text{ giờ } 12\text{ phút} = 72\text{ phút}$.\\
Quãng đường chạy được là:
$$250 \times 72 = 18\,000\text{ (m)} = 18\text{ km}.$$
\textbf{Đáp án của câu hỏi là B}
}
\end{ex}

\begin{ex} %[5D10V60]
Một con thuyền đi xuôi dòng từ A đến B hết $2$ giờ $30$ phút. Biết vận tốc của thuyền khi nước lặng là $18\text{ km/h}$ và vận tốc dòng nước là $2\text{ km/h}$. Quãng đường sông AB dài là
\choice
{$45\text{ km}$}
{\True $50\text{ km}$}
{$52\text{ km}$}
{$40\text{ km}$}
\loigiai{
Vận tốc của thuyền khi xuôi dòng là:
$$18 + 2 = 20\text{ (km/h)}.$$
Đổi: $2\text{ giờ } 30\text{ phút} = 2{,}5\text{ giờ}$.\\
Quãng đường AB dài là:
$$20 \times 2{,}5 = 50\text{ (km)}.$$
\textbf{Đáp án của câu hỏi là B}
}
\end{ex}

\begin{ex} %[5D10V60]
Một người đi từ nhà đến cơ quan mất $24$ phút nếu đi với vận tốc $12\text{ km/h}$. Nếu hôm nay người đó muốn đến cơ quan sớm hơn $6$ phút thì phải đi với vận tốc bao nhiêu?
\choice
{$14\text{ km/h}$}
{\True $16\text{ km/h}$}
{$15\text{ km/h}$}
{$18\text{ km/h}$}
\loigiai{
Đổi: $24\text{ phút} = 0{,}4\text{ giờ}$.\\
Quãng đường từ nhà đến cơ quan là:
$$12 \times 0{,}4 = 4{,}8\text{ (km)}.$$
Thời gian đi để đến sớm hơn $6$ phút là:
$$24 - 6 = 18\text{ (phút)} = \dfrac{18}{60}\text{ giờ} = 0{,}3\text{ giờ}.$$
Vận tốc cần đi là:
$$4{,}8 : 0{,}3 = 16\text{ (km/h)}.$$
\textbf{Đáp án của câu hỏi là B}
}
\end{ex}

\begin{ex} %[5D10V60]
Một vận động viên đua xe đạp đi $3$ vòng quanh một bờ hồ. Vòng thứ nhất đi hết $15$ phút, vòng thứ hai đi với vận tốc $24\text{ km/h}$ hết $12$ phút, vòng thứ ba đi hết $13$ phút. Biết mỗi vòng bờ hồ dài $4{,}8\text{ km}$. Tổng quãng đường và thời gian vận động viên đó hoàn thành $3$ vòng đua là
\choice
{$14{,}4\text{ km}$ và $45\text{ phút}$}
{\True $14{,}4\text{ km}$ và $40\text{ phút}$}
{$15\text{ km}$ và $40\text{ phút}$}
{$14{,}4\text{ km}$ và $38\text{ phút}$}
\loigiai{
Tổng quãng đường đi $3$ vòng là:
$$4{,}8 \times 3 = 14{,}4\text{ (km)}.$$
Tổng thời gian đi là:
$$15 + 12 + 13 = 40\text{ (phút)}.$$
\textbf{Đáp án của câu hỏi là B}
}
\end{ex}

\begin{ex} %[5D10V60]
Khoảng cách giữa hai thành phố là $240\text{ km}$. Một ô tô xuất phát từ thành phố thứ nhất lúc $8\text{ giờ } 15\text{ phút}$ với vận tốc $60\text{ km/h}$. Giữa đường ô tô dừng nghỉ $30$ phút. Ô tô đến thành phố thứ hai lúc
\choice
{$12\text{ giờ } 15\text{ phút}$}
{\True $12\text{ giờ } 45\text{ phút}$}
{$13\text{ giờ}$}
{$13\text{ giờ } 15\text{ phút}$}
\loigiai{
Thời gian ô tô chạy trên đường là:
$$240 : 60 = 4\text{ (giờ)}.$$
Tổng thời gian đi và nghỉ là:
$$4\text{ giờ} + 30\text{ phút} = 4\text{ giờ } 30\text{ phút}.$$
Thời điểm ô tô đến nơi là:
$$8\text{ giờ } 15\text{ phút} + 4\text{ giờ } 30\text{ phút} = 12\text{ giờ } 45\text{ phút}.$$
\textbf{Đáp án của câu hỏi là B}
}
\end{ex}

\begin{ex} %[5D10V60]
Một máy bay cất cánh từ TP. Hồ Chí Minh lúc $10\text{ giờ } 20\text{ phút}$ và hạ cánh xuống Hà Nội lúc $12\text{ giờ } 30\text{ phút}$. Biết vận tốc bay trung bình của máy bay là $800\text{ km/h}$. Chiều dài đường bay từ TP. Hồ Chí Minh đến Hà Nội xấp xỉ là
\choice
{$1\,680\text{ km}$}
{\True $1\,733\text{ km}$}
{$1\,600\text{ km}$}
{$1\,800\text{ km}$}
\loigiai{
Thời gian bay là:
$$12\text{ giờ } 30\text{ phút} - 10\text{ giờ } 20\text{ phút} = 2\text{ giờ } 10\text{ phút} = \dfrac{13}{6}\text{ (giờ)}.$$
Độ dài quãng đường bay là:
$$800 \times \dfrac{13}{6} \approx 1\,733{,}33\text{ (km)} \approx 1\,733\text{ km}.$$
\textbf{Đáp án của câu hỏi là B}
}
\end{ex}

\begin{ex} %[5D10V60]
Một con rùa bò với vận tốc $8\text{ cm/giây}$. Để bò hết quãng đường $36\text{ m}$, con rùa cần thời gian là
\choice
{$400\text{ giây}$}
{\True $450\text{ giây}$}
{$500\text{ giây}$}
{$360\text{ giây}$}
\loigiai{
Đổi: $36\text{ m} = 3\,600\text{ cm}$.\\
Thời gian con rùa bò hết quãng đường là:
$$3\,600 : 8 = 450\text{ (giây)} = 7\text{ phút } 30\text{ giây}.$$
\textbf{Đáp án của câu hỏi là B}
}
\end{ex}

\textbf{Phần 2. Trắc nghiệm đúng sai}
\setcounterpageref{ex}{0}

\begin{ex} %[5D10V60]
Một người đi xe máy khởi hành từ A lúc $7\text{ giờ}$ sáng với vận tốc $36\text{ km/h}$ để đến B cách A một khoảng $90\text{ km}$. Xét tính đúng/sai của các nhận định sau:
\choiceTF
{\True Thời gian người đó đi từ A đến B (nếu không nghỉ) là $2{,}5$ giờ}
{\True Nếu không nghỉ dọc đường, người đó đến B lúc $9\text{ giờ } 30\text{ phút}$}
{Nếu dọc đường người đó dừng nghỉ $45$ phút thì đến B lúc $10\text{ giờ}$}
{\True Nếu muốn đến B lúc $9\text{ giờ}$ sáng và không nghỉ thì người đó phải đi với vận tốc $45\text{ km/h}$}
\loigiai{
\begin{itemchoice}
\item Đúng. $t = 90 : 36 = 2{,}5\text{ giờ} = 2\text{ giờ } 30\text{ phút}$.
\item Đúng. $7\text{ giờ} + 2\text{ giờ } 30\text{ phút} = 9\text{ giờ } 30\text{ phút}$.
\item Sai. Nếu nghỉ 45 phút thì đến B lúc $9\text{ giờ } 30\text{ phút} + 45\text{ phút} = 10\text{ giờ } 15\text{ phút}$ (không phải $10$ giờ).
\item Đúng. Từ 7 giờ đến 9 giờ là 2 giờ. Vận tốc cần đi là $90 : 2 = 45\text{ km/h}$.
\end{itemchoice}
}
\end{ex}

\begin{ex} %[5D10V60]
Hai bạn Lan và Mai cùng xuất phát từ trường đi đến thư viện huyện trên quãng đường dài $6\text{ km}$. Lan đi bộ với vận tốc $4\text{ km/h}$, Mai đi xe đạp với vận tốc $12\text{ km/h}$. Xét tính đúng/sai của các nhận định sau:
\choiceTF
{\True Mai đi hết quãng đường trong thời gian $30$ phút}
{\True Lan đi hết quãng đường trong thời gian $1$ giờ $30$ phút}
{Mai đến thư viện trước Lan $45$ phút}
{\True Khi Mai đến nơi thì Lan đã đi được quãng đường dài $2\text{ km}$}
\loigiai{
\begin{itemchoice}
\item Đúng. $t_{Mai} = 6 : 12 = 0{,}5\text{ giờ} = 30\text{ phút}$.
\item Đúng. $t_{Lan} = 6 : 4 = 1{,}5\text{ giờ} = 1\text{ giờ } 30\text{ phút} = 90\text{ phút}$.
\item Sai. Mai đến trước Lan: $90 - 30 = 60\text{ phút} = 1\text{ giờ}$ (chứ không phải $45$ phút).
\item Đúng. Khi Mai đến nơi (sau $0{,}5$ giờ), quãng đường Lan đã đi được là $4 \times 0{,}5 = 2\text{ (km)}$.
\end{itemchoice}
}
\end{ex}

\textbf{Phần 3. Câu trả lời ngắn}
\setcounterpageref{ex}{0}

\begin{ex} %[5D10V60]
Một ô tô đi với vận tốc $54\text{ km/h}$ trong $2$ giờ $20$ phút. Quãng đường ô tô đi được là bao nhiêu ki-lô-mét?
\shortans{126}
\loigiai{
Đổi: $2\text{ giờ } 20\text{ phút} = \dfrac{7}{3}\text{ giờ}$.\\
Quãng đường ô tô đi được là:
$$54 \times \dfrac{7}{3} = 126\text{ (km)}.$$
}
\end{ex}

\begin{ex} %[5D10V60]
Một người đi bộ với vận tốc $5\text{ km/h}$ trên quãng đường dài $12\text{ km}$. Hỏi người đó đi hết bao nhiêu phút?
\shortans{144}
\loigiai{
Thời gian đi hết quãng đường là:
$$12 : 5 = 2{,}4\text{ (giờ)}.$$
Đổi ra phút:
$$2{,}4 \times 60 = 144\text{ (phút)}.$$
}
\end{ex}

\begin{ex} %[5D10V60]
Một người đi xe máy từ A đến B với vận tốc $36\text{ km/h}$. Sau khi đi được $1$ giờ $45$ phút thì người đó còn cách B một đoạn dài $17\text{ km}$. Hỏi toàn bộ quãng đường AB dài bao nhiêu ki-lô-mét?
\shortans{80}
\loigiai{
Đổi: $1\text{ giờ } 45\text{ phút} = 1{,}75\text{ giờ}$.\\
Quãng đường đã đi được là:
$$36 \times 1{,}75 = 63\text{ (km)}.$$
Độ dài quãng đường AB là:
$$63 + 17 = 80\text{ (km)}.$$
}
\end{ex}

\begin{ex} %[5D10V60]
Một ca nô đi ngược dòng trên khúc sông dài $42\text{ km}$ với vận tốc $14\text{ km/h}$. Hỏi ca nô đi hết khúc sông đó trong bao nhiêu giờ?
\shortans{3}
\loigiai{
Thời gian ca nô đi hết khúc sông là:
$$42 : 14 = 3\text{ (giờ)}.$$
}
\end{ex}

\begin{ex} %[5D10V60]
Một xe đạp điện chạy với vận tốc $25\text{ km/h}$. Để đi hết quãng đường $35\text{ km}$, xe đạp điện cần thời gian là bao nhiêu phút?
\shortans{84}
\loigiai{
Thời gian đi quãng đường đó là:
$$35 : 25 = 1{,}4\text{ (giờ)}.$$
Đổi ra phút:
$$1{,}4 \times 60 = 84\text{ (phút)}.$$
}
\end{ex}

\begin{ex} %[5D10V60]
Một người chạy xe máy từ thị trấn về làng cách nhau $60\text{ km}$. Biết vận tốc xe là $40\text{ km/h}$. Người đó muốn về tới làng lúc $11\text{ giờ } 30\text{ phút}$ trưa thì phải xuất phát muộn nhất lúc mấy giờ? (Nhập số giờ nguyên, ví dụ 10 giờ thì nhập 10).
\shortans{10}
\loigiai{
Thời gian đi từ thị trấn về làng là:
$$60 : 40 = 1{,}5\text{ (giờ)} = 1\text{ giờ } 30\text{ phút}.$$
Thời điểm xuất phát muộn nhất là:
$$11\text{ giờ } 30\text{ phút} - 1\text{ giờ } 30\text{ phút} = 10\text{ giờ}.$$
}
\end{ex}
